% New English translation from checked cote 25, site 84 (PDF page 85).
% AI-translated by GPT-6 Astra (OpenAI, ChatGPT Pro mode).
\scanpage{25/84}
\mspage{\hand{78}}

14. \uline{Conical projections}.

NB We have already used conical projections on various
\var{\margin{$\|$}}{\placedmargin{$\|$}}{marginal note placed}{ega25_25_084-v0001}occasions, in particular at the end of no. 8, the formulation of 10.4 and elsewhere, and
the following sorites should probably come instead at the beginning of the paragraph,
and possibly in the auxiliary “Grassma\unsure{n}nian” paragraph.
Let $C = P(F)$ be a linear subvariety of $P(E)=P$, of relative
dimension $r-m-1$ \var{\xout{in}}{}{typewriter strike-out}{ega25_25_084-v0002} over $S$\var{\inserted[typed]{i.e. of codimension $m+1$ in $P$}}{ i.e. of codimension $m+1$ in $P$}{insertion}{ega25_25_084-v0003}, so $F$ is a quotient of $E$, loc free of
rank $r-m$, $F=E/G$, where $G$ is loc free of rank $m+1$. We defined
algebraically in Chap II a morphism
\[ p_C : P - C = P(E) - P(E/G) \hand{\longrightarrow} P(G) , \]
which we shall interpret geometrically, and which will be called
(in view of the following description) the \uline{conical projection} with center $C$.
(NB We assume $r-m-1$ lies between $-1$ and $r-1$, i.e. $m$ between $0$
and $r$, with no further restriction.) For this, let us begin by interpreting \var{\xout{\unsure{Q(G)}}}{}{typewriter strike-out}{ega25_25_084-v0004} $P(G)$
as a sub\var{\inserted[typed]{pre}}{pre}{insertion}{ega25_25_084-v0005}scheme\var{\inserted[typed]{closed}}{ closed}{insertion}{ega25_25_084-v0006} in $\mathrm{Grass}^m(P)= \mathrm{Grass}_{r-m+1}(P)$, by means of the evident homomorphism
of functors
\[ \var{\xout{\unsure{Gr|G}}}{}{typewriter strike-out}{ega25_25_084-v0007}\ P(G) \hand{\longrightarrow} \mathrm{Grass}_{r-m+1}(E) \]
obtained by considering, for each invertible quotient $G/G'$ of $G$, the
locally free Module of rank $(r-m)+1$ $E/G'$ of $E$ (and the same after
any base change). The preceding homomorphism of functors is a
\uline{monomorphism}, and since the first is proper over $S$ and the second separated,
it is a closed immersion. \var{\xout{\unsure{(}}}{}{typewriter strike-out}{ega25_25_084-v0008}More generally, we should make explicit
the closed immersions of the Grassmannians of $G$, i.e. of $P(G)$,
into those of \var{\inserted{$E$ i.e. of}}{$E$ i.e. of}{insertion}{ega25_25_084-v0009} $P=P(E)$. \var{\xout{\unsure{Here,}}}{}{typewriter strike-out}{ega25_25_084-v0010} The image\var{\inserted[typed]{(in the functor sense)}}{ (in the functor sense)}{insertion}{ega25_25_084-v0011} of the resulting morphism consists of
the linear subvarieties \var{\xout{of M}}{}{typewriter strike-out}{ega25_25_084-v0012} $L$\var{\inserted[typed]{of the desired dimension}}{ of the desired dimension}{insertion}{ega25_25_084-v0013} of $P$ that \uline{contain} $C$. Denoting
this image by $Q(C)$ in the case at hand (i.e. for the dimensions
specified above), and identifying $P(G)$ with $Q(C)$, the conical
projection morphism
\[ p_C : P - C \hand{\longrightarrow} Q(C) \hand{\subset} \mathrm{Grass}^m(P) \]
\tnote{25-084-1}{Conical projection renders projection conique. The subscript r-m+1 in the displayed Grassmannians and the source's description of E/G' as a Module of E are retained. Kif kif is rendered the same.}

\mspage{\hand{84}}
